\section{Real and complex analysis from integrable presentations}
\label{sec:analytic-apps}

The preceding constructions support continuous mathematics whose input
need not have pointwise values everywhere. We give two applications with
complete quantitative proofs: a contractive integral equation and a
holomorphic transform of an \(L^1\) name. They also illustrate why an
operator must be measured in a norm appropriate to its use.

\subsection{The operator norm needed by an integral equation}

For a rational rectangular step kernel \(K(x,t)\) on \([0,1]^2\), define
\[
 (T_Ku)(x)=\int_0^1 K(x,t)u(t)\,dt,\qquad
 \|K\|_{\mathrm{col}}=\sup_t\int_0^1 |K(x,t)|\,dx.
\]
These expressions are finite sums and a finite maximum. On common
dyadic levels \(a\) and \(b\), with coefficients \(k_{ij}\),
\[
 \|K\|_{\mathrm{col}}=
    \max_{j<2^b}2^{-a}\sum_{i<2^a}|k_{ij}|.
\]
Refinement leaves this value unchanged. The finite triangle inequality
and interchange of the two finite sums give
\begin{equation}
 \|T_Ku\|_1\leq\|K\|_{\mathrm{col}}\|u\|_1,
 \qquad
 \|(T_K-T_L)u\|_1\leq\|K-L\|_{\mathrm{col}}\|u\|_1.
 \label{eq:colbound}
\end{equation}
Thus a regular name of kernels in the column norm, with a uniform norm
bound, defines a bounded operator on the represented \(L^1\) space.
The finite-sum proof is the object-theory kernel; extension uses only
the indicated precision shifts.

The product \(L^1\) norm of a kernel would not suffice here. For
\(K_m(x,t)=2^m\mathbf1_{[0,2^{-m})}(t)\), its product \(L^1\) norm is
one but its column norm is \(2^m\). Applied to
\(u_m(t)=2^m\mathbf1_{[0,2^{-m})}(t)\), a unit \(L^1\) input, the
output is the constant \(2^m\). All these quantities are finite rational
calculations. This is an operational reason to record observation norms
in the framework.

\subsection{A primitive-recursive resolvent construction}

\begin{theorem}[Contractive integral equation]\label{thm:resolvent}
Suppose \(K_s\) and \(g_s\) are regular kernel and \(L^1\) names,
respectively, with finite certificates
\[
 \|K_s\|_{\mathrm{col}}\leq q<1,\quad
 \|g_s\|_1\leq M,\quad
 \|K_s-K_t\|_{\mathrm{col}}\leq2^{-s}+2^{-t},\quad
 \|g_s-g_t\|_1\leq2^{-s}+2^{-t}.
\]
Here \(q,M\) are rational, \(0\leq q<1\). Choose an integer \(b\geq2\)
with \(q\leq1-1/b\), and put
\[
 C=b+Mb^2,\quad s(k)=k+h(4C),\quad
 R(k)=k+h(4Mb)+1,\quad N(k)=bR(k).
\]
The finite step functions
\[
 y(k)=\sum_{j=0}^{N(k)-1}T_{K_{s(k)}}^j g_{s(k)}
 \label{eq:resolvent-name}\tag{RN*}
\]
form a regular \(L^1\) name. Its represented value is the unique solution
of \(f=g+T_Kf\), and \(\|f\|_1\leq M/(1-q)\).
\end{theorem}

\begin{proof}
We first prove finite estimates. For any operator \(A\) with
\(\|Au\|_1\leq q\|u\|_1\), induction on \(j\) gives
\(\|A^ju\|_1\leq q^j\|u\|_1\). Hence the norm of a partial Neumann
sum is at most \(M\sum_{j<n}q^j\leq Mb\), and any tail beginning at
\(n\) is at most \(Mbq^n\). All sums are finite at this step.

For two contractions \(A,B\), the finite identity
\[
 A^j-B^j=\sum_{i=0}^{j-1}A^{j-1-i}(A-B)B^i
\]
is proved by expanding the sum and canceling neighboring terms. Applied
to a vector of norm at most \(M\), it yields the bound
\(j q^{j-1}M\|A-B\|\) when \(j>0\). A direct finite geometric-sum
calculation gives
\(\sum_{j=1}^{n-1}j q^{j-1}\leq(1-q)^{-2}\leq b^2\).
Consequently two partial sums of the same length, with inputs \(u,v\)
and operators \(A,B\), differ by at most
\[
 b\|u-v\|_1+Mb^2\|A-B\|.
 \tag{P}\label{eq:resolvent-perturb}
\]
Comparing two lengths with their common maximum adds at most the two
tail bounds.

To bound the number of iterations without logarithmic real analysis,
Bernoulli's finite inequality gives
\((1+1/(b-1))^{b-1}\geq2\). Thus
\((1-1/b)^b\leq1/2\), and \(q^{bR}\leq2^{-R}\).
It follows that
\(C2^{-s(k)}\leq2^{-k}/4\) and
\(Mbq^{N(k)}\leq2^{-k}/8\). Combining these bounds with
\eqref{eq:resolvent-perturb} shows
\[
 \|y(k)-y(\ell)\|_1
 \leq\tfrac38(2^{-k}+2^{-\ell}),
\]
which is stronger than regularity.

At the finite level,
\[
 y(k)-g_{s(k)}-T_{K_{s(k)}}y(k)
       =-T_{K_{s(k)}}^{N(k)}g_{s(k)}.
\]
Its norm is at most \(Mq^{N(k)}\). Replacing the input and kernel by
any finer named approximations costs at most
\((1+Mb)(2^{-s(k)}+2^{-t})\), by \eqref{eq:colbound}.
These rational bounds prove the equation as equality of names.
If two solutions differ by \(w\), then \(w=T_K^nw\) for every finite
\(n\), so a rational bound \(D\) on its norm gives
\(\|w\|_1\leq Dq^n\). Choose \(n\) by the same explicit geometric
schedule to make this less than any requested dyadic tolerance.
The names are equal. The partial-sum norm bound gives the stated norm
of the solution.
\end{proof}

The theorem proves a substantial fragment of Fredholm-type integral
equations. It requires a certified contraction gap, not merely a
classical assertion of invertibility. In particular, it does not
extract a gap from a spectral-radius assertion or choose a solution at
a singular parameter.

\paragraph{Worked example.}
Take \(K(x,t)=xt/2\) and \(g(x)=1\). The column norm is \(1/4\).
If \(f(x)=1+cx\), then
\[
 (T_Kf)(x)=x\left(\tfrac14+\tfrac c6\right).
\]
The equation is solved by \(c=3/10\), hence
\(f(x)=1+3x/10\). This identity is checked by finite polynomial
integration. Dyadic step approximations to \(K\) converge in the
column norm with an explicit mesh bound, so the construction above
produces the same name without using an exact ansatz. The example
also checks the orientation of the column-norm bound.

\subsection{A holomorphic transform of rough input}

Let \(f\) be a real \(L^1([0,1])\) name with
\(\|f\|_1\leq M\). There need be no effective pointwise evaluator for
\(f\). For \(|z|<1\), define the represented transform
\[
 G_f(z)=\int_0^1\frac{f(t)}{1-zt}\,dt.
 \tag{G}\label{eq:analytic-transform}
\]
The integral means extension of a bounded linear functional on step
functions, not evaluation of an unspecified measurable representative.
For rational complex inputs, a guard \(|z|\leq r<1\) is certified by
\((\Re z)^2+(\Im z)^2\leq r^2\). No square-root oracle is needed.
More generally, a rational upper bound \(|a+ib|\leq Q\), with
\(Q\geq0\), is the squared rational inequality \(a^2+b^2\leq Q^2\).
The triangle and product estimates used below can therefore be certified
by finite rational inequalities, without rationality of the modulus
itself.

\begin{theorem}[Analytic observation of an \(L^1\) name]
Fix rational \(0\leq r<1\), put \(\delta=1-r\), and take an integer
\(b\geq2\) with \(\delta^{-1}\leq b\). On the closed disk
\(|z|\leq r\), the transform is constructively defined and satisfies
\[
 |G_f(z)-G_g(z)|\leq b\|f-g\|_1,
 \qquad
 |G_f(z)-G_f(w)|\leq Mb^2|z-w|.
 \label{eq:g-stability}\tag{GS}
\]
It is complex differentiable, with
\[
 D_f(z)=\int_0^1\frac{t f(t)}{(1-zt)^2}\,dt,
 \qquad
 |D_f(z)-D_g(z)|\leq b^2\|f-g\|_1.
\]
Whenever \(|z|,|z+h|\leq r\),
\begin{equation}
 |G_f(z+h)-G_f(z)-hD_f(z)|\leq Mb^3|h|^2.
 \label{eq:g-derivative}
\end{equation}
Its Taylor coefficients at zero are the moments
\(a_j=\int_0^1t^jf(t)\,dt\), and
\begin{equation}
 \left|G_f(z)-\sum_{j=0}^{N-1}a_jz^j\right|
       \leq Mb r^N.
 \label{eq:g-taylor}
\end{equation}
All approximations and schedules can be chosen primitive recursively
from the input presentation and the stated rational guards.
\end{theorem}

\begin{proof}
The denominator has modulus at least \(1-|z|t\geq\delta\).
The finite identities
\[
 \frac1{1-zt}-\frac1{1-wt}
       =\frac{(z-w)t}{(1-zt)(1-wt)},
\]
and
\[
 \frac1{1-(z+h)t}-\frac1{1-zt}
        -\frac{ht}{(1-zt)^2}
       =\frac{h^2t^2}{(1-(z+h)t)(1-zt)^2}
\]
give the bounds after multiplication by \(|f|\) and integration.
Every estimate is first a weighted finite-sum inequality for step
data and then passes to names by a uniform operator bound. The
quadratic remainder proves differentiability directly; it does not
invoke a general theorem about interchanging limits and integrals.
At a point on the boundary of the displayed disk, use the wider rational
guard \(r'=(1+r)/2<1\) when taking arbitrary sufficiently small
increments \(h\). The same argument then proves the full derivative
test there. This wider guard also accommodates rational approximations
to a represented complex input: choose their error below \(r'-r\)
before evaluation, and use the \(z\)-Lipschitz estimate to account for
that error. The sharper bounds stated for pairs inside the original
\(r\)-disk remain unchanged.

For explicit production, take a step approximant \(f_s\) and refine
its partition to mesh \(2^{-m}\). Midpoint quadrature of
\(f_s(t)/(1-zt)\) is a complex rational sum. The kernel is
\(b^2\)-Lipschitz in \(t\), so its quadrature error is at most
\(\|f_s\|_1b^2 2^{-m}\). Input error is at most \(b2^{-s}\)
when the name has been interpreted with its certified limit error.
Choose both budgets below \(2^{-k-2}\), and require \(m\) at least
the level of the step partition. In the entirely finite comparison
version, replace \(2^{-s}\) by \(2^{-s}+2^{-t}\); the same shifts
give a regular name.

For the derivative kernel, the elementary difference estimate on
\([0,1]\) is
\[
 \left|\frac{t}{(1-zt)^2}-\frac{u}{(1-zu)^2}\right|
 \leq (b^2+4b^4)|t-u|\leq5b^4|t-u|.
\]
Indeed the first numerator change costs \(b^2|t-u|\); the change of
the squared denominator costs at most
\(|z||t-u|\,|2-z(t+u)|b^4\leq4b^4|t-u|\).
Midpoint quadrature therefore supplies the derivative name with a
primitive-recursive mesh shift as well.

Finally the finite geometric identity
\[
 \frac1{1-zt}=\sum_{j<N}(zt)^j+\frac{(zt)^N}{1-zt}
\]
proves \eqref{eq:g-taylor}. For a step function constant at \(c_i\)
on \([u_i,v_i)\), its moments are exactly
\[
 a_j(f_s)=\sum_i c_i
       \frac{v_i^{j+1}-u_i^{j+1}}{j+1}.
 \tag{M}\label{eq:moments}
\]
The coefficient error is at most the input \(L^1\) error, uniformly in
\(j\). A partial polynomial therefore costs at most \(b\) times that
input error. The tail is made dyadically small by
\(N=b(k+h(4Mb)+1)\), using the elementary geometric bound from the
previous proof. Thus one step approximant and one table of exact
moments suffice for all coefficients up to \(N\). The entire proof
uses finite identities and explicit remainder bounds.
\end{proof}

For the density reconstructed from the dyadic masses of
\(f(t)=2(1+t)/3\), the moments are
\[
 a_j=\frac23\left(\frac1{j+1}+\frac1{j+2}\right).
\]
This gives an exact rational test family for the moment code. It also
illustrates the intended connection: finite mass data first constructs
an integrable object, after which stable integral observations construct
an analytic object. No pointwise limit, general Radon--Nikodym oracle,
or unrestricted dominated-convergence principle is inserted between
the stages.
